% Extended supplement methods — moved out of the journal supplement (reviewer point 9) to reduce
% its page count. This is repository material accompanying the manuscript
% "A Comparability Protocol for Benchmarking Relational Database Access Layers in Express.js"
% (Zenodo 10.5281/zenodo.21313858 concept DOI). It carries the reporting checklist, the
% benchmarking-pitfalls checklist, the detailed statistical methods, and the measurement details.
\documentclass[11pt,a4paper]{article}
\usepackage[utf8]{inputenc}
\usepackage[T1]{fontenc}
\usepackage{booktabs}
\usepackage{array}
\usepackage{amsmath}
\usepackage{enumitem}
\usepackage[protrusion=true,expansion=false]{microtype}
\usepackage{url}
\title{Extended supplement: reporting checklist, pitfalls checklist,\\
  statistical methods, and measurement details\\[2pt]
  \large A Comparability Protocol for Benchmarking Relational Database Access Layers in Express.js}
\author{Mateusz Miotk\\ University of Gda\'nsk}
\date{}
\begin{document}
\maketitle
\noindent This document holds material moved out of the numbered online supplement to keep the
journal supplement compact; it is distributed with the replication package. The numbered supplement
tables and figures (Supplement Tables~S1--S44, Figures~S1--S4) remain in \texttt{supplement.pdf}.

\section{A reporting checklist for access-layer benchmarks}
\label{sup:checklist}
Distilled from the design decisions and the pitfalls this study confronted, the
following items are the ones an access-layer benchmark should report so that its
claims can be scoped and reproduced. It is a reporting aid, not a standard.
\begin{enumerate}
  \item \textbf{Correctness before timing.} Assert byte-identical responses across
        layers (final serialized bodies), or state what differs; several defects in
        this study were caught only this way.
  \item \textbf{Experimental unit.} Treat the freshly booted run, not the individual
        request, as the unit; requests within a run are not independent.
  \item \textbf{Same-host repetition is not environmental replication.} Report how
        many runs, on how many hosts, over what period; do not call same-host
        repeats ``independent.''
  \item \textbf{Documentation-selected treatment.} State the exact API, its documentation page and
        version, the declined alternative, and whether logging, hooks, validation,
        and identity maps are on.
  \item \textbf{Latency construct.} Distinguish high-load closed-loop response time
        (which tracks throughput and queueing) from utilization-controlled or
        open-loop latency; report the operating point.
  \item \textbf{Same-SQL / shared-plan control.} Report it as a compound intervention
        (it changes query formulation, protocol, preparation, round-trips, and
        mapping together), not a mechanism decomposition.
  \item \textbf{Resource accounting.} Report application-tier and database CPU
        separately; throughput per CPU-second, not throughput alone.
  \item \textbf{Write-state control.} Reset the mutable state between runs and state
        the durability regime (default versus relaxed).
  \item \textbf{Provenance.} Pin engine images by digest and dependencies by
        lockfile; archive raw per-cell records, seeds, and a table-to-record manifest.
  \item \textbf{Scope.} State the host topology, workload, versions, and dates, and
        confine claims to them.
\end{enumerate}


\section{Benchmarking-pitfalls checklist}
\label{sup:pitfalls}
The four confounds summarized in the main text's Discussion, in full. Each is a generic failure
mode of the genre, easy to introduce by accident and hard to detect: the correctness ones
without an automated cross-check, the performance ones without profiling every cell.
\begin{itemize}
  \item \textbf{Aggregation fan-out.} A naive join between \texttt{posts} and \texttt{comments}
        evaluated before the \texttt{SUM} multiplies each post's view count by its number of
        matching comment rows, inflating the aggregate --- a correctness bug invisible unless
        every layer's output is checked against a reference, not a performance one.
  \item \textbf{OFFSET vs.\ keyset pagination.} An early, OFFSET-based design for the range-scan
        endpoint collapsed on MySQL at realistic offsets, which would have been misread as an
        access-layer weakness rather than what it is: a pagination-strategy cost that InnoDB
        pays more dearly for than PostgreSQL. Rewriting every adapter around keyset pagination
        (\texttt{WHERE id < cursor ORDER BY id DESC LIMIT n}) removed the confound.
  \item \textbf{Write-workload contamination.} The insert endpoint grows \texttt{posts} on
        every call; because engines and layers commit at different rates, later cells in a run
        would scan a differently sized table depending on run order rather than on layer
        identity, unless benchmark-inserted rows are deleted before every cell.
  \item \textbf{Query-formulation confounds.} A first attempt at the aggregation endpoint
        pre-aggregated \emph{all} of \texttt{comments} before filtering to one author, so every
        request re-scanned the whole table regardless of which author was requested --- a
        query-plan artifact masquerading as an access-layer effect, fixed by rewriting the query
        as a correlated subquery touching only the target author's rows.
\end{itemize}

\section{Statistical methods}
\label{sup:stats}
This section details the tests summarized in the main text's Analysis subsection. In every
test the unit is the \emph{run-level} statistic (one throughput or p99 value per repeated
run, 25 per cell), never the individual request; all confidence intervals are percentile
bootstraps resampling replicate indices (preserving the pairing) from a fixed seed
(\texttt{mulberry32(0x57a75)}), so every resample is reproducible.

To test whether adjacently ranked layers are distinguishable, we compare each adjacent pair
on the deep/nested fetch on their per-replicate throughput ratios with a two-sided paired
sign-flip permutation test ($20{,}000$ permutations, so the smallest attainable $p$ is
$1/(20{,}001)$), cross-checked with the Wilcoxon signed-rank test, and report the
geometric-mean paired ratio with a paired bootstrap 95\% CI and the paired dominance (the
fraction of replicates the faster layer wins). Because p99 is reported at millisecond
resolution and contains ties, the permutation test keeps zero log-ratios and the Wilcoxon
test drops them with the standard tie correction. Because the seven adjacent pairs follow
the \emph{observed} ranking rather than a preregistered one, their significances are a
secondary, descriptive robustness check, and we apply a Bonferroni correction as a
conservatism check rather than claiming a confirmatory family.

The layer$\times$engine interaction is tested per pattern with a blocked permutation test:
for each layer and replicate we form the paired PostgreSQL-minus-MySQL log-throughput
difference $D_{\ell,i}$ and, under the null of no interaction, permute the layer labels of
$\{D_{\ell,i}\}$ \emph{within} each replicate $i$ ($20{,}000$ permutations); the statistic
is the between-layer $F$ of a randomized-block ANOVA on $D$ with replicate as the block.
Rank agreement across engines (Spearman, Kendall) is reported descriptively over the seven
evaluated layers --- systems under study, not a sample from a population --- so we attach no
population confidence intervals to those correlations.

For the \texttt{pg}--Prisma pair we additionally report a paired two-one-sided-tests (TOST)
equivalence result against a $\pm5\%$ margin on the log-ratio. That margin is an
author-selected smallest effect of practical interest for \emph{this} study, not a general
engineering threshold: at modest deployment scale a sub-$5\%$ throughput difference would
not change a capacity-planning decision (for example, how many application instances to
provision for a target load). We set it from that engineering consequence rather than from
measurement noise; it happens also to sit within the run-to-run and day-to-day variation
typical of a virtualized host (this study's own run-to-run CV reaches $15.9\%$ on the
noisiest MySQL-insert cell, within $4.2\%$ on the reads), but that coincidence is not the
justification. At very large scale even $5\%$ can correspond to many machines, so the margin
should be reset per deployment. We chose it before the analysis but did not preregister it,
and the equivalence conclusion is not sensitive to the exact value within a few percent.

\section{Measurement details}
\label{sup:measure}
This section records the measurement specifics summarized in the main text's Measurement
procedure.

\paragraph{Warm-up length}
The fifteen-second warm-up follows a separate cold-boot measurement for a fast, a middling,
and the slowest PostgreSQL layer: every layer reached and sustained at least 95\% of
steady-state throughput within 12~s (Prisma by 5~s, native by 9~s, MikroORM by 12~s),
absorbing JIT compilation, pool fill, and plan caching.

\paragraph{Order invariance and completeness}
Cell order was reshuffled independently within each replicate; a forward-versus-reversed
order check on the write endpoint (5 replicates per direction) found no detectable history
effect on the read cells (reversed-to-forward ratio 0.977--1.280, median 1.012; the wide
tail is confined to the noisy MySQL inserts, where five replicates cannot separate order
from write noise). Of the 450 dedicated write-endpoint boots, 449 completed normally; the
exception was one Drizzle/MySQL write replicate whose server boot did not clear the
post-reset health check, leaving that cell with 24 of 25 runs (the runner records and skips
such failures rather than retrying silently).

\paragraph{Tail estimation}
\texttt{autocannon} records a per-run latency histogram at millisecond resolution; even the
slowest deep-fetch cells ($\approx480$--$516$~req/s over the 12~s window) complete
$\approx5{,}800$--$6{,}200$ requests per run, so each run-level p99 is estimated from
roughly 60 requests in its upper $1\%$ tail, and the reported p99 is the median of the 25
run-level values.

\paragraph{Resource and framework sampling}
Server CPU and peak resident memory were sampled from \texttt{/proc} over the full process
tree (parent plus descendants, so an in-process engine or spawned helper cannot escape
measurement); a \texttt{perf\_hooks} GC observer (event count, total pause time) and the
200~ms connection-pool sampler ran inside the server, polled through a \texttt{/stats}
endpoint after each run. A fixed-payload \texttt{/baseline} endpoint returning a
seed-realistic thread-shaped payload measured the shared Express-plus-serialization floor:
10{,}909 (PostgreSQL server) and 10{,}921 (MySQL server) req/s, roughly $3\times$ the
fastest deep-fetch cell, so the framework floor leaves between-layer spreads visible rather
than compressed.

\paragraph{Insert configuration}
Each insert is an autocommit single-statement \texttt{INSERT} under each engine's default
isolation level (PostgreSQL Read Committed, MySQL Repeatable Read); the transactional
variant wraps a post and its five comments in one transaction. The default durability
settings are PostgreSQL \texttt{fsync=on}, \texttt{synchronous\_commit=on},
\texttt{full\_page\_writes=on} and MySQL \texttt{innodb\_flush\_log\_at\_trx\_commit=1},
\texttt{sync\_binlog=1}, \texttt{log\_bin=ON}.


\end{document}
